\subsection{Preparing for Confession}\label{sec:confession}

The homily the Office reads on the Fourth Sunday goes on from the day's Gospel to John's words to
the crowds that came to be baptized: who warned you to flee from the wrath to come? St Gregory
answers:

\begin{quote}
``The wrath to come'' in one sense signifieth the great vengeance of the Latter Day: the sinner that
repenteth not of his sin now, will have no mean whereby to flee from punishment then.\par
{\smallskip\noindent\small ---Gregory the Great, \textit{Homiliae in Evangelia} 20, tr. Bute,
\textit{The Roman Breviary} (1908), Winter, p.~259.\par}
\end{quote}

\noindent The season that looks for the Judge also offers the remedy before he comes. St Francis de
Sales, writing to his Philothea, puts it simply: ``Our Saviour has left the holy
sacrament of penance and confession to his Church, that in it we might cleanse ourselves from all
our iniquities, as often as we should be defiled by them. Never suffer your heart then, O
Philothea! to remain long infected with sin, since you have so easy a remedy at hand''
(\textit{Introduction to a Devout Life} II.19).

\subsubsection{What the Church teaches}

The Council of Trent, in its fourteenth session, names three acts of the penitent: ``the acts of
the penitent himself, to wit, contrition, confession, and satisfaction, are, as it were, the matter
of this sacrament,'' and they ``are for this reason called the parts of penance.'' Of the three it
puts contrition first:

\begin{quote}
Contrition, which holds the first place amongst the aforesaid acts of the penitent, is a sorrow of
mind, and a detestation for sin committed, with the purpose of not sinning for the future.\par
{\smallskip\noindent\small ---Council of Trent, Session XIV, chapter~4, tr. J. Waterworth,
\textit{The Canons and Decrees of the Council of Trent}, p.~95.\par}
\end{quote}

\noindent The Baltimore Catechism gives the same teaching in a single line: ``The substance of an Act of
Contrition is: I am sorry for my sins, because they have offended God, and I will never sin again''
(\textit{Catechism of Christian Doctrine} No.~3, Q.~108). Trent also names the fruit of the
sacrament: its effect ``is reconciliation with God, which sometimes, in persons who are pious and
who receive this sacrament with devotion, is wont to be followed by peace and serenity of conscience,
with exceeding consolation of spirit'' (Session XIV, chapter~3, tr. Waterworth, p.~95).

\subsubsection{An examination by the four Sundays}

The book offers, as its own counsel, an examination drawn from the texts of the season. It does not
replace the examination of the commandments that a catechism or prayer book provides; it adds the
season's questions to it. Take one part a night in the days before confession.

\begin{enumerate}[leftmargin=1.5em, itemsep=0.5ex]
\item \textit{The First Sunday: the works of darkness.} Romans 13 names rioting and drunkenness,
chambering and wantonness, strife and envying. Where in the past months have I made provision for
the flesh? Have I slept when I should have been awake to God, in prayer neglected or put off?
\item \textit{The Second Sunday: the reed.} Gregory reads the reed shaken by the wind as the mind
that bends with praise and blame. Have I done good to be seen, or left good undone for fear of
what others would say? Have I been ashamed of Christ, or of his Cross, before others?
\item \textit{The Third Sunday: the truth about oneself.} John confessed what he was and what he
was not. Have I claimed what is not mine, in word or in my own mind? Have I lied, flattered,
boasted, or let a false impression stand? Have I refused a duty that was mine?
\item \textit{The Ember Days: the handmaid.} Mary called herself the handmaid of the Lord and went
in haste to serve. Have I resisted God's will where I knew it? Have I been slow to help those near
me who needed it?
\item \textit{The Fourth Sunday: the crooked and the rough.} Gregory's list for the generation of
vipers is envy of the good, evil repaid with evil, and ways sought to harm a neighbor. Is there
injustice in my dealings that I have not set right? Is there harshness in my temper, at home or at
work, that I have excused?
\end{enumerate}

After the examination, make an act of contrition. The text printed in
Appendix~\ref{sec:prayers} is the one the Baltimore \textit{Manual of Prayers} gave the Catholics
of the United States in 1889.

\subsubsection{When to go}

The rule of life (Section~\ref{sec:rule}) suggests confession once before Christmas at the least,
and in Ember week for those who can. Go early: a confession made in the second or third week leaves the
rest of Advent to live out its purpose of amendment. St Francis de Sales counselled his Philothea to go much more often:
``Confess yourself humbly and devoutly once every week, and always, if possible, before you
communicate, although your conscience should not reproach you with the guilt of mortal sin''
(\textit{Introduction to a Devout Life} II.19). That is his counsel to a soul seeking devotion, and
a confessor can say how far it suits another. His reason holds for anyone: ``for by confession you
not only receive absolution from the venial sins you confess, but likewise strength to avoid them,
light to discern them well, and grace to repair all the damage you may have sustained by them.''

\medskip
{\noindent\small St Francis de Sales is quoted from the English \textit{Introduction to a Devout
Life} published by Pustet (New York, undated; translator not named).\par}
